% =============================================================
% IV. VALIDATION, LIMITATIONS AND SENSITIVITY
% =============================================================
\section{Validation, limitations and sensitivity}
\label{sec:limitations}

\subsection{The overlaps}

The reconstruction can be tested at two points. The interwar series, built from Estonian sources and weighted by the 1938 benchmark, and the Soviet-period series, built from Soviet and archival sources and weighted by the 1990 benchmark, both cover the years 1937--40; and the Soviet-period series and the modern series, backcast from the MPD's 2011 benchmark, both cover 1990--5. Table~\ref{tab:overlaps} reports the results.

% tab:overlaps
\begin{table}[htbp]
\centering\small
\caption[Overlapping estimates of Estonian GDP per capita at the two splicing points]{\textsc{Overlapping estimates of Estonian GDP per capita at the two splicing points}}
\label{tab:overlaps}
\begin{threeparttable}
\begin{tabular}{lrrcc}
\toprule
Year & Earlier-benchmark series & Later-benchmark series & Difference (\%) & Feinstein grade \\
& (\intlnineteen{}) & (\intlnineteen{}) & & \\
\midrule
\multicolumn{5}{l}{\textit{(a) Interwar series (1938 benchmark) and Soviet-period series (1990 benchmark)}} \\
1937 & 2,722 & 2,678 & $-$1.6 & Firm \\
1938 & 2,685 & 2,782 & +3.6 & Firm \\
1939 & 2,736 & 2,654 & $-$3.0 & Firm \\
1940$^{\dagger}$ & 2,776 & 2,767 & $-$0.3 & Firm \\
\midrule
\multicolumn{5}{l}{\textit{(b) Soviet-period series (1990 benchmark) and adjusted MPD series (2011 benchmark)}} \\
1990 & 8,425 & 8,172 & +3.0 & Firm \\
1991$^{\dagger}$ & 7,506 & 7,558 & $-$0.7 & Firm \\
1992 & 5,400 & 6,069 & $-$12.4 & Good \\
1993 & 4,916 & 5,870 & $-$19.4 & Rough \\
1994 & 4,973 & 5,899 & $-$18.6 & Rough \\
1995 & 4,989 & 6,277 & $-$25.8 & Conjectural \\
\bottomrule
\end{tabular}
\begin{tablenotes}[para,flushleft]
\footnotesize
\textit{Sources:} see online appendix; adjusted MPD series from \cite{mpd2023} with 1990--2010 growth replaced by World Bank constant-price growth backcast from the 2011 benchmark (\cite{wdi2024}). \textit{Notes:} Difference = (later $-$ earlier)/earlier in panel (a) and (earlier $-$ later)/earlier in panel (b), so that a positive number always means that the 1990-based series is the higher one. Feinstein's grades: firm ${<}5$ per cent, good 5--15, rough 15--25, conjectural ${>}25$ (\cite[21]{feinstein1972}). $^{\dagger}$~Splice year: the final series takes the arithmetic average of the two estimates.
\end{tablenotes}
\end{threeparttable}
\end{table}


In 1937--40 the two series differ by $-$1.6, $+$3.6, $-$3.0 and $-$0.3 per cent. All four differences fall within the 5 per cent margin that Feinstein reserved for `firm' estimates, and they do not share a sign.\footnote{\autocite[21]{feinstein1972}.} The test is a demanding one. The two series share no sources: the interwar indices come from the republic's own statistics and the Soviet-period indices from the yearbooks of the ESSR and the post-war archives. They share no weights, since one uses the sectoral structure of 1938 and the other that of 1990. And they are anchored to benchmarks constructed by different methods and separated by fifty-two years, one a bilateral price comparison with Sweden, the other a multilateral purchasing-power-parity study of the Soviet republics. That the Soviet-period series, extrapolated backwards over half a century of war, collectivisation and forced industrialisation, arrives within 4 per cent of an independently measured level is the strongest evidence available that the net error of the extrapolation is small. It cannot exclude errors that happen to offset, such as an overstatement of growth from hidden inflation in trade turnover against an understatement from the index-number assumptions. But the two index-number biases discussed below, late base-year weights and fixed value-added weights, would both understate growth and so would both place the backcast 1937--40 level above the interwar series; the overlap therefore tests their combined size. A net understatement of growth of as little as 0.1 percentage point a year over 1940--90 would place the Soviet-period series 5 per cent above the interwar series in every overlap year, which is not observed. Broadberry and Burhop proposed that time-series projections should be regarded as consistent with direct benchmarks when they agree within 10 per cent; a projection over thirty years in their own work landed within 4 per cent of the benchmark and was judged a success.\footnote{\autocite[323--6]{broad2007}; \autocite[933]{broad2008}.} The Estonian overlap meets that standard over a longer span, and does so in each of four consecutive years.

The second overlap is less tidy, and the reasons are instructive. In 1990 the Soviet-period series stands 3.0 per cent above the modern series and in 1991 0.7 per cent below it; thereafter the gap widens to 12 per cent in 1992 and 26 per cent by 1995. The close agreement in 1990--1 is itself a validation of the 1990 benchmark, and in particular of the conversion from roubles to dollars on which a reviewer might fasten. Bolotin's dollar figure for Estonia was obtained with purchasing-power parities, not exchange rates, and the modern series against which it is tested derives from an entirely independent ICP round, that of 2011, carried backwards with the World Bank's constant-price growth rates. Two routes that share no data arrive at the same level of Estonian income in 1990--1 within 3 per cent. The divergence after 1991 arises because the Soviet-period series continues to value the physical collapse of Soviet-era output at 1990 weights and, in trade, to rely on deflated turnover figures that understate inflation: it records what happened to the Soviet-era economy, not what replaced it. The modern series is therefore preferred from 1991, the year in which the two agree, and the MPD's own growth rates for 1990--2010, which are chained to its NMP-based level for the late Soviet years, are replaced by the World Bank's. The MPD's 1989 figure for Estonia lies 32 per cent above the new series (Table~\ref{tab:previous}); the implied fall in income during the transition in the MPD is 36 per cent, against \eeTROUGHpct{} per cent here, so that the two sources agree on the depth of the collapse even though they disagree on the level from which it started.\footnote{\autocite{mpd2023}; the comparison is developed in \autocite{klimantas2026soviet}.}

\subsection{A single late benchmark and the index-number problem}

The Soviet-period series extrapolates 1990 weights backwards over 1937--90, a period of profound structural change. An index with a late base year assigns to each sector the relative prices of the end of the period, and since the sectors that grew fastest were, as a rule, those whose relative prices fell, late weights give less importance to fast-growing sectors and so understate growth; conversely, early weights overstate it. Feinstein and Thomas put the point generally: an index with an early base `will always show a greater increase in quantities than one constructed with a base year late in the period'.\footnote{\autocite[516]{feinsteinthomas2002}.} This is the effect that Gerschenkron first documented for Soviet machinery, and it was large in the Stalin period: Bergson's GNP for 1928--37 grows by 175 per cent at 1928 prices and by 62 per cent at 1937 prices.\footnote{\autocite[453]{gerschenkron1957}; \autocite[517]{feinsteinthomas2002}.} Three considerations establish that it is far smaller for the series presented here.

First, the effect is concentrated in the early phases of industrialisation, when relative prices change fastest; Jonas and Sardy showed that `most of the action which leads to the Gerschenkron effect takes place early'.\footnote{\autocite[85]{jonassardy1970}.} For the post-war Soviet economy the measured effect is modest. When the CIA moved its base year from 1970 to 1982 prices, its estimate of Soviet GNP growth over 1951--84 fell by about half a percentage point a year, from 5.6 to 5.2 per cent for the 1950s and from 3.2 to 2.4 per cent for the 1970s; an earlier comparison of 1960, 1970 and 1976 prices produced differences of a few tenths of a point, which the CIA judged `not large'.\footnote{\autocite[8]{cia1990}; \autocite[23--4]{cia1982a}; \autocite[156]{harrison1993}.} Half a point a year over forty years compounds to roughly a fifth of the level; this indicates the order of magnitude to be expected for the post-war decades rather than a bound, since the Estonian extrapolation also spans the forced industrialisation of the 1940s and 1950s, in which the effect is largest. Second, the bound must therefore come from the data, and the series can be tested. If 1990 weights understated Estonian growth over 1940--90, the Soviet-period series would arrive in 1937--40 systematically above the interwar series, which is weighted by the prices of 1938. It does not: the four differences are small and of both signs. The structural change within industry, where the Gerschenkron effect originates, is moreover handled at a finer level than the aggregate weights suggest, since the 101 product series are aggregated within eleven branches with 1986 weights before the branches are combined. Third, the procedure is the one used by every reconstruction of Soviet output with which the series might be compared. The CIA used 1970 weights for 1950--80 and 1982 weights for 1950--90, and Maddison, reviewing the practice, found `nothing untoward' in it; Kuboniwa's and Ponomarenko's reconstructions of Russian GDP for 1961--90 use, among their variants, the fixed structure of Russia's 1990 GDP.\footnote{\autocite[312]{Maddison1998b}; \autocite[57, 62]{ponomarenko2002}.} A chain-linked alternative would require annual current-price value-added by sector, which does not exist for any Soviet republic, and a 1938-weighted alternative for the whole period would require the sectors of 1990 to be valued at prices of 1938, when several of them did not exist. The reader who wishes to see the question settled empirically should turn to the 1937--40 overlap.

\subsection{Fixed value-added weights}

A related assumption is that the ratio of value-added to gross output within each sector remained at its benchmark-year level. The volume indices track gross output, or inputs, or employment, and are combined with value-added weights; the method is sometimes called single deflation, although the term properly describes the deflation of nominal value-added by an output price index, a procedure that this paper does not use and that performs much worse.\footnote{\autocite[443--4]{thomasfeinstein2003}.} Thomas and Feinstein's evaluation of the method actually used here, output indicators with value-added weights, found that it understated true value-added growth by between 0.15 and 1 percentage point a year in their simulations, depending on the pace of technical change, and they described it as `the next best procedure' to double deflation where intermediate inputs cannot be observed.\footnote{\autocite[427--8, 444]{thomasfeinstein2003}.} Double deflation is not available for either period. The interwar statistical office recorded no intermediate consumption, and for the Soviet period, as Ponomarenko notes, the price statistics that double deflation would require `practically did not exist'; the CIA applied double deflation only to agriculture and used gross output as the index of value-added everywhere else, on the argument, due to Hill, that gross output is a better proxy than inputs unless output and input growth diverge by more than two points a year.\footnote{\autocite[18--19]{ponomarenko2002}; \autocite[45]{cia1982a}.}

The bias is bounded with the help of the Swedish historical national accounts, which permit single- and double-deflated series to be compared over comparable phases of industrialisation.\footnote{\autocite[19]{swedengdp_15602010}.} For the interwar period, when input costs outside agriculture moved little, the maximum error is about 2 per cent of the level of GDP. For the Soviet period the issue is more serious, because the sustained shift towards intermediate-intensive heavy industry makes it likely that value-added ratios changed, and the maximum error is put at 14 per cent of the level of GDP over the half-century, the derivation being given in the online appendix. The sign adopted in Table~\ref{tab:sensitivity} is that of Thomas and Feinstein's simulations, an understatement of growth and hence an overstatement of the backcast early levels; 14 per cent over fifty years corresponds to about a quarter of a percentage point a year, in the lower part of their range, which is appropriate because the Estonian indices track gross output product by product rather than by broad branch. Two observations limit the practical concern. The error accumulates over the extrapolation, so that it is largest in the earliest years of the Soviet-period series, precisely where the 1937--40 overlap shows the series to be within 4 per cent of an independent level. And the opposite sign cannot be excluded, since the Soviet shift towards intermediate-intensive heavy industry would have raised rather than lowered value-added ratios if the rising degree of fabrication outpaced the relative fall in intermediate-input prices; this is why the overlap, not the simulations, is decisive. Gerschenkron himself, confronting constant gross-to-net ratios over half a century in Jostock's German accounts, thought them `not so very implausible', since the rising degree of fabrication is likely to have been offset by the relative fall in the prices of value-added components.\footnote{\autocite[453]{gerschenkron1957}.}

\subsection{Other sources of bias}

Table~\ref{tab:sensitivity} lists the remaining sources of bias, the sectors they affect and the maximum impact of each on the level of GDP; the bounds for the occupation period are derived in the same way as in the companion Lithuanian study, with the Estonian sectoral weights.

% tab:sensitivity
\begin{table}[htbp]
\centering\footnotesize\setlength{\tabcolsep}{4pt}
\caption[Sources of potential bias and their maximum impact on the level of GDP]{\textsc{Sources of potential bias and their maximum impact on the level of GDP}}
\label{tab:sensitivity}
\begin{threeparttable}
\begin{tabular}{p{2.5cm}p{6.6cm}p{2.4cm}r}
\toprule
Period & Source of bias & Sectors affected & Max. impact \\
& & & (\% of GDP) \\
\midrule
Interwar, 1919--40 & Service-sector value-added in the 1938 benchmark & Services (34\% of GDP) & $\pm$3 \\
 & Fixed value-added weights (`single deflation') & All & $\pm$2 \\
 & Industrial output proxied by employment and real wages & Industry & $\pm$1 \\
 & Taxes less subsidies not tracked annually & All & $\pm$1 \\
 & No productivity adjustment in public services & Public services & $\pm$3 \\
 & Other services proxied by population & Other services & $\pm$1 \\
 & Military personnel known only for 1919--20 and 1939 & Defence & $\pm$1 \\
\midrule
Soviet period, 1940--90 & Service-sector allocation in the 1990 benchmark & Services & $\pm$4 \\
 & Taxes less subsidies absorbed in sectoral values & All & $-$0.5 \\
 & Fixed value-added weights (`single deflation') & All & $-$14 \\
 & Hidden inflation in comparable-price trade turnover & Trade & +2 \\
 & Size of Soviet forces stationed in Estonia unknown & Defence & $\pm$5 \\
 & No productivity adjustment in employment-based services & Services & $-$10 \\
 & Wartime movements 1943--5 partly borrowed from Lithuania & All (1941--5 only) & 15--25 \\
\bottomrule
\end{tabular}
\begin{tablenotes}[para,flushleft]
\footnotesize
\textit{Sources:} authors' assessment; see Section~\ref{sec:limitations} and the online appendix for the derivation of each bound. \textit{Notes:} signed entries indicate the direction in which the published series would be revised if the bias were corrected; $\pm$ entries are symmetric. The bounds are not additive.
\end{tablenotes}
\end{threeparttable}
\end{table}


Three deserve comment. The service sector of the 1938 benchmark, about one-third of GDP, rests on earnings and employment data and on Latvian structural ratios; an error of 10 per cent in the service estimate moves GDP by 3 per cent, and Norkus and Markevi\v{c}i\={u}t\.{e}'s independent approximation, 5 per cent below the new benchmark, suggests that any error is of that order. In the 1990 benchmark the allocation of value-added among services depends on Bolotin's Union-wide parities and the 1987 input--output table; a reallocation of a quarter of service value-added between branches moves GDP by about 4 per cent. The Soviet-period trade index uses turnover at comparable prices, the one value series in the reconstruction; Kontorovich's estimates of hidden inflation in Soviet retail trade imply an upward bias of at most 2 per cent of GDP.\footnote{\autocite[688--9]{kontorovich_2001}.} Defence after 1940 is unobserved, since the Soviet forces stationed in Estonia were an all-Union institution and their size is unknown between 1940 and 1990; the defence component of public services is interpolated between the 1940 and 1990 totals in step with civil employment, and the maximum bias is 5 per cent of GDP in either direction. The employment-based indices for health, education and administration make no allowance for productivity change, which understates service growth if productivity rose; the maximum effect is 10 per cent of the level. Finally, unregistered activity is estimated to have amounted to 10--20 per cent of Soviet GDP in 1989, but OECD economies displayed shadow economies of comparable size at the same date, so that the omission does not systematically impair international comparison.\footnote{\autocite[532--51]{kim2003}; \autocite[5]{baltictransitioncontraction}; \autocite[90--1]{shadow_europe}.}

Taken together, the quantifiable biases imply a margin of about 6 per cent around the interwar levels and, around the Soviet-period levels, of 14 per cent for the largest single bias, or about a fifth in the earliest years if the two same-signed biases of Table~\ref{tab:sensitivity}, fixed value-added weights and unadjusted service productivity, were both at their maxima; the overlap shows that the realised error is far smaller. The war years are subject to a wider margin. The sub-period growth rates on which the argument of the following sections rests are robust to all of them: no plausible combination of the biases listed in Table~\ref{tab:sensitivity} would raise growth in 1968--89 to the level of 1950--68, close the gap with Finland in 1938, or reverse Estonia's position relative to the Soviet average.
